\normalsize % Mengembalikan font ke 12pt untuk isi artikel

\input{jurnal/konten-utama/c1-pendahuluan}
\section{TINJAUAN PUSTAKA}

\subsection{Sub Judul 1}
\PlaceholderParShort

\subsection{Sub Judul 2}
\PlaceholderParShort
\input{jurnal/konten-utama/c3-metode-penelitian}
\section{HASIL DAN PEMBAHASAN}

\subsection{Sub Judul 1}
\PlaceholderParShort

\subsection{Sub Judul 2}
\PlaceholderParShort
\section{PENUTUP}

\subsection{Simpulan}
\PlaceholderPar

\subsection{Saran}
\PlaceholderPar

% =======================================================
% OTOMATISASI DAFTAR PUSTAKA MENGGUNAKAN APA STYLE
% =======================================================
\section{REFERENSI}
% 1. Menghapus teks bawaan daftar pustaka (karena kita sudah buat \section{REFERENSI})
\renewcommand{\bibsection}{}
% 2. Mengubah ukuran font khusus daftar pustaka menjadi 11pt sesuai aturan
\small
% 3. Mengatur hanging indent sebesar 1.25cm (menjorok ke dalam)
\setlength{\bibhang}{1.25cm}
% 4. Mengatur jarak antar entri daftar pustaka (8pt)
\setlength{\bibsep}{8pt}

\nocite{*} % Memanggil semua entri yang ada di file .bib
\bibliography{a7-pustaka-digunakan}